\documentclass[12pt]{article}
\usepackage{amsmath, amssymb, geometry, enumitem}
\geometry{margin=0.65in}
\setlength{\parindent}{0pt}
\setlength{\parskip}{0.1em}

\title{ECON 101: Macroeconomics \\ Solutions -- Quiz 5}
\author{Instructor: Muhebullah Karimzada}
\date{November 26,  2025}

\begin{document}
\maketitle

\section*{Part A --- Multiple Choice Solutions}

\setlist[enumerate]{itemsep=3pt, topsep=3pt}

\textbf{1. A decrease in Canadian interest rates will:}\\
\textbf{Correct answer: (B) Shift AD right.}

\textbf{Reasoning:}
\begin{itemize}
    \item Lower interest rates reduce the cost of borrowing.
    \item This stimulates investment spending by firms and durable consumption by households.
    \item In the AD--AS model, higher planned spending at each price level shifts the \emph{aggregate demand curve} to the right.
    \item It is not a movement along AD (price level has not changed) and does not directly shift SRAS.
\end{itemize}

\bigskip

\textbf{2. Which of the following would shift the SRAS curve to the left?}\\
\textbf{Correct answer: (C) An increase in corporate taxes.}

\textbf{Reasoning:}
\begin{itemize}
    \item Short-run aggregate supply depends on input costs (wages, raw materials), productivity, and factors that affect profitability.
    \item An \emph{increase in corporate taxes} raises firms' costs and reduces after-tax profits.
    \item Firms are less willing to supply output at each price level, so SRAS shifts \emph{left}.
    \item A decrease in wages and a fall in oil prices reduce costs and would shift SRAS right.
    \item A technological improvement increases productivity and also shifts SRAS right.
\end{itemize}

\newpage

\textbf{3. The long-run aggregate supply (LRAS) curve:}\\
\textbf{Correct answer: (C) Is vertical at potential GDP.}

\textbf{Reasoning:}
\begin{itemize}
    \item In the long run, output is determined by \emph{real factors}: labour, capital, technology, and institutions.
    \item The price level does not affect the full-employment level of output.
    \item Therefore the LRAS curve is vertical at potential GDP (\(Y^*\)).
    \item It does not slope upward; that describes SRAS.
\end{itemize}

\bigskip

\textbf{4. The money multiplier equals:}\\
\textbf{Correct answer: (B) $\dfrac{1}{r_d}$.}

\textbf{Reasoning:}
\begin{itemize}
    \item In the simple deposit expansion model, the \emph{simple deposit multiplier} is:
    \[
    m = \frac{1}{r_d},
    \]
    where \(r_d\) is the desired reserve ratio.
    \item The other options either invert or modify this expression incorrectly.
\end{itemize}

\bigskip

\textbf{5. If the Bank of Canada sells government securities in the open market, the most likely outcome is:}\\
\textbf{Correct answer: (B) Higher interest rates and lower AD.}

\textbf{Reasoning:}
\begin{itemize}
    \item When the Bank of Canada \emph{sells} bonds, commercial banks and the public pay for these bonds.
    \item This \emph{reduces} bank reserves and the money supply.
    \item With a smaller supply of money, interest rates tend to rise.
    \item Higher interest rates discourage investment and some consumption, decreasing aggregate demand.
\end{itemize}
\newpage

\section*{Part B --- Numerical / Analytical Problems}

\textbf{Problem 1 (7 marks) --- AD--AS \& Short-Run Equilibrium}

\textit{Given:}
\[
Y^* = 2{,}500 \text{ (billion dollars)}, \quad P = 120, \quad \Delta I = -50 \text{ (billion)}, \quad MPC = 0.75.
\]
Assume the price level is fixed for the multiplier calculation.

\subsection*{(a) Compute the expenditure multiplier. (1 mark)}

The simple spending (Keynesian) multiplier is:
\[
k = \frac{1}{1 - MPC}.
\]

Substitute \(MPC = 0.75\):
\[
k = \frac{1}{1 - 0.75} = \frac{1}{0.25} = \boxed{4}.
\]

\subsection*{(b) Calculate the total change in real GDP. (2 marks)}

Initial change in autonomous spending:
\[
\Delta A = \Delta I = -50 \text{ (billion)}.
\]

Total change in equilibrium real GDP:
\[
\Delta Y = k \cdot \Delta A = 4 \times (-50) = \boxed{-200 \text{ (billion)}}.
\]

\textit{Interpretation:} A \$50 billion fall in investment results in a \$200 billion fall in equilibrium GDP.

\subsection*{(c) What is the new level of real GDP? (1 mark)}

Initial real GDP was at potential:
\[
Y_0 = Y^* = 2{,}500.
\]

After the shock:
\[
Y_1 = Y_0 + \Delta Y = 2{,}500 + (-200) = \boxed{2{,}300 \text{ (billion)}}.
\]

\subsection*{(d) AD--AS short-run effects on GDP, unemployment, and price level. (3 marks)}

\textbf{Step-by-step explanation:}
\begin{itemize}
    \item The reduction in investment shifts the \emph{aggregate demand} curve leftward (from AD\(_0\) to AD\(_1\)).
    \item In the short run, SRAS is upward-sloping and LRAS is vertical at \(Y^* = 2{,}500\).
    \item The new short-run equilibrium is at the intersection of AD\(_1\) and SRAS.
    \item As a result:
    \begin{itemize}
        \item \textbf{Real GDP} falls from \(2{,}500\) to \(2{,}300\): the economy is in a \emph{recessionary gap}.
        \item \textbf{Unemployment} rises because firms reduce output and hire fewer workers.
        \item \textbf{Price level:} In the full AD--AS model, the fall in demand also puts downward pressure on the price level. 
        \item In this question, we used a fixed price level for the multiplier calculation, but conceptually, in the AD--AS diagram, both real GDP and the price level fall in the short run.
    \end{itemize}
\end{itemize}


\vspace{3em}

\textbf{Problem 2 (8 marks) --- Simple Deposit Multiplier}

\textit{Given:}
\[
\text{Initial new deposit} = \$5{,}000, \quad r_d = 0.10.
\]

\subsection*{(a) Compute the simple deposit multiplier. (1 mark)}

Formula:
\[
m = \frac{1}{r_d}.
\]

Substitute \(r_d = 0.10\):
\[
m = \frac{1}{0.10} = \boxed{10}.
\]

\subsection*{(b) Maximum possible increase in total deposits. (2 marks)}

The maximum total change in deposits is:
\[
\Delta D_{\max} = m \times \text{initial new deposit}.
\]

\[
\Delta D_{\max} = 10 \times 5{,}000 = \boxed{\$50{,}000}.
\]

\textit{Interpretation:} If banks lend out all excess reserves and all loans are redeposited in the banking system, total deposits can rise by \$50,000.

\subsection*{(c) Required reserves and loans from the initial deposit. (3 marks)}

\textbf{Step 1: Compute desired reserves from initial deposit.}
\[
\text{Desired reserves} = r_d \times \text{initial deposit} = 0.10 \times 5{,}000 = \boxed{\$500}.
\]

\textbf{Step 2: Compute excess reserves (amount that can be loaned out).}
\[
\text{Excess reserves} = \text{initial deposit} - \text{desired reserves} 
= 5{,}000 - 500 = \boxed{\$4{,}500}.
\]

\textit{Interpretation:}  
From the initial \$5,000:
\begin{itemize}
    \item The bank must keep \$500 as reserves.
    \item It can make new loans of up to \$4,500, which then get redeposited and re-lent through the system, generating the full deposit expansion (up to \$50,000 in total deposits in the simple model).
\end{itemize}

\subsection*{(d) Why the actual increase may be smaller than the maximum. (2 marks)}

\textbf{Any one of the following (with explanation) is acceptable:}
\begin{itemize}
    \item \textbf{Cash leakage:} Some of the loan proceeds may be held by the public as cash rather than redeposited, breaking the multiplier chain.
    \item \textbf{Excess reserves:} Banks may choose to hold excess reserves above the required/desired level, so they do not lend out the full amount of excess reserves.
    \item \textbf{Credit risk / weak loan demand:} In a recession or uncertain environment, banks may be reluctant to lend or borrowers may be reluctant to take on new loans.
\end{itemize}

For full marks, the student should state clearly that these behaviours reduce the effective multiplier, so the actual increase in deposits and the money supply is \emph{less} than the simple deposit multiplier predicts.

\hrule
\vspace{1em}

\begin{center}
\textit{End of Solutions}
\end{center}

\end{document}
